\section{Introduction}
\label{sec:introduction}

The Smith--Toda realization problem asks whether the first chromatic
parameters can be killed in a finite spectrum. Algebraically these
parameters form a regular sequence in Brown--Peterson homology.
Topologically, killing the next parameter requires an actual map of
finite spectra that induces its multiplication on homology. Regularity
computes the homology of the cofiber once such a map exists; it does
not supply the map. At a prime $p$, write
\[
 \BP_* = \Z_{(p)}[v_1,v_2,\ldots],
 \qquad |v_i|=2(p^i-1),
\]
using Hazewinkel generators. A Smith--Toda complex $V(n)$ is a finite
$p$-local spectrum $X$ with
\[
 \BP_*X\cong\BP_*/(p,v_1,\ldots,v_n)
\]
as graded $\BP_*\BP$-comodules, with the quotient coaction and its
generator in degree zero. Here finite $p$-local means a retract of the
$p$-localization of a finite CW spectrum. The question concerns first
powers. Suitable higher powers can be realized at every finite height
by the periodicity methods of Hopkins--Smith
\cite[Theorem~9 and Proposition~5.14]{HopkinsSmith1998}. Those
exponents are chosen as part of the construction; the question here
prescribes every exponent to be one.

\begin{theorem}\label{thm:main}
For every integer $n\geq0$, there exist a prime $p$ and a finite
$p$-local spectrum $X$ such that
\[
 \BP_*X\cong\BP_*/(p,v_1,\ldots,v_n)
\]
as graded $\BP_*\BP$-comodules, with the canonical quotient coaction
and generator in degree zero.
\end{theorem}

Thus the all-height Smith--Toda existence problem has a positive answer
when the prime may depend on the height. The construction takes place
in ordinary $p$-local spectra, before any chromatic localization. It
does not assert that one prime works at every height, provide an
explicit bound on the prime, or require a multiplication on the final
spectrum.

\subsection{Context and earlier methods}
The realization problem grew out of the study of complex cobordism
modules and periodic phenomena in stable homotopy. Milnor's work on
complex cobordism and Novikov's coefficient theorem established its
polynomial coefficient structure \cite{Milnor1960,Novikov1962}.
Brown and Peterson introduced the spectrum $\BP$ \cite{BrownPeterson1966}.
Quillen
identified the formal group law underlying that theory and its
$p$-typical projection \cite{Quillen1969}; Hazewinkel supplied
the integral generators and coefficient formulas used here
\cite{HazewinkelIII}. These results describe the algebraic quotients
whose finite topological realization is at issue.

The low-height constructions are due to Adams, Smith, and Toda:
$V(k)$ exists for $k=1,2,3$ when $p>2k$, respectively, as summarized
in \cite[p.~493]{Nave2010}. Adams's periodicity construction is
part of his work on the groups $J(X)$ \cite{Adams1966}.
Smith's work treats the realization
of complex bordism modules \cite{Smith1970}, while Toda constructs
spectra realizing exterior parts of the Steenrod algebra
\cite{Toda1971}. At height zero the Moore spectrum supplies the
quotient by $p$. Each subsequent first-power attachment asks for
additional topological information. In the September 2023 version of
\cite[Remark~3.29]{CulverZhang2024}, the existence of $V(4)$ was
recorded as open at every prime. We address arbitrary fixed height
by allowing the prime to grow.

Known nonexistence results are consistent with this quantifier order.
For $p\geq7$, Nave proves that $V((p+1)/2)$ does not exist and that
$V((p-1)/2)$ cannot be a ring spectrum if it exists
\cite[Theorem~1.3]{Nave2010}. The obstructed height depends on $p$,
and the second assertion concerns additional structure. Neither is an
obstruction at one fixed height at every prime.

Ultraproducts exploit the distinction between a fixed height and a
varying prime. Barthel--Schlank--Stapleton establish asymptotic
algebraicity in chromatically localized categories and apply it to
generalized Moore objects \cite{BSS2020}. Their compactness is in
the localized category, whereas Theorem~\ref{thm:main} requires
an ordinary finite spectrum. Here the ultraproduct is formed from
ordinary $p$-local spectra. We retain varying sphere degrees and
construct a finite sequence of cofibers there. Only finitely many
maps and homotopies must subsequently descend to an individual
prime.

The other ingredients have complementary predecessors. Exactness
of the fixed-prime Tate tensor power provides a model for the
permutation construction \cite[Proposition~III.1.1]{NS2018}; our
comparison requires a single Euler exponent for the entire family
of tensor cubes as $p$ varies. The normalization uses completion
by a perfect algebra, in the setting of the Amitsur methods of
Mathew--Naumann--Noel \cite{MNN2017}. Power operations in cobordism
go back to tom Dieck and Quillen \cite{tomDieck1968,Quillen1971}.
For detection, we use the coherent Thom construction and its
Euler-product formula \cite{May1977,JohnsonNoel2010}. We retain these powers in complex
cobordism throughout the calculation, then apply the ordinary
multiplicative projection to Brown--Peterson homology.

A related use of filtrations appears in Randal-Williams's construction
of Smith--Toda analogues in graded algebra modules
\cite[Section~5, Theorem~5.6 and Lemma~5.7]{RandalWilliams2026}.
There successive cofibers are controlled through associated-graded
generators and vanishing estimates over a field of fixed positive
characteristic. The comparison here is methodological: our
connectivity bounds vary with the prime and feed the normalization
of an Euler layer in ordinary spectra. Recent local realization
results of Kato--Shimomura--Shimomura concern quotients with a
chromatic generator inverted \cite[Theorem~1.7]{KatoShimomura2025}.
They address a different realization target from the ordinary
finite, first-power quotient considered here.

\subsection{Construction and Detection}
Fix $n\geq1$ and a nonprincipal ultrafilter $\U$ on the odd primes.
The central ambient category is
\[
 \cC=\prod_{\U}\Sp_{(p)}.
\]
Its mapping spectra are rational, although sequences of spheres in
degrees growing with $p$ still carry essential information. A formal
weight $w_j$ records the varying even degree $2(p^j-1)$. Mapping out
of these spheres gives a rational graded algebra in which finite
Koszul extensions can be formed without discarding that degree data.

In this overview, $S^{2(p^j-1)}$ denotes the sequence of the indicated
primewise spheres, and $\one$ denotes the tensor unit. We construct
objects $Y_0,\ldots,Y_{n+1}$ in $\cC$, starting from $Y_0=\one$,
and, for $0\leq j\leq n$, maps
\[
 b_j:S^{2(p^j-1)}\longrightarrow Y_j,
 \qquad
 Y_{j+1}=\cofib\bigl(
 S^{2(p^j-1)}\otimes Y_j\xrightarrow{\ b_j\cdot-\ }Y_j
 \bigr).
\]
The first class $b_0$ is represented by primewise multiplication by
$p$. The multiplication used in this display
exists in $\cC$; the proof does not require its full commutative
structure to exist at any individual prime.

There are two separate tasks. The first is to construct $b_{j+1}$
from $b_j$, including the choices of boundaries that permit iteration.
The second is to show that these classes kill the intended first
powers in homology. Keeping the tasks separate avoids assuming the
homology calculation in order to construct the next class.

For the construction, take the $p$-fold permutation power with the
affine group $C_p\rtimes\F_p^{\times}$ acting on the factors. The
main difficulty is uniformity: the diagrams comparing powers with
cofibers have dimensions that grow with $p$. We prove that one Euler
map, associated to twice the reduced permutation representation,
annihilates the entire comparison object. In a suitable auxiliary
category, inverting that map consequently makes the power exact.
Before inversion its Euler parameter $t$ has a quotient that detects
the next actual sphere degree. The argument controls the whole
comparison object, rather than merely its individual filtration pieces.

The power must be normalized before it can be iterated on a Koszul
cycle. We do this on a universal algebra $R$ of cycle coefficients,
with a finite-cell Koszul extension $K$. Unlike the constructed $Y_j$,
these universal coefficient algebras have actual primewise commutative
models to which the power applies. At weights proportional to
$p$, a symmetric-group connectivity estimate gives the necessary
vanishing only after tensoring over $R$ with $K$. Such a tensor test
need not detect equivalences on arbitrary modules. We show that it
does detect the comparisons needed by our actual completed targets,
using the Amitsur totalization associated to $R\to K$. This permits
normalization by powers of $t$ while preserving algebra maps and
chosen nullhomotopies, not just operations on homotopy classes.
For the attaching cycle, it removes a prescribed Euler factor while
retaining a class before inversion, where reduction modulo $t$ is
still available.

For detection, let $x_j$ be the coefficient of $Z^{p^j}$ in the
$p$-series of complex cobordism, with $x_0=p$. A primewise
Thom-corrected power calculation identifies the leading $t$-term
of the powered $x_j$ with a unit times $x_{j+1}$. Comparison with
the normalized construction proves that the Hurewicz image of
$b_j$ is a unit times $x_j$ modulo its predecessors. Under the
ordinary multiplicative projection to $\BP$, the ideals generated
by the $x_j$ become the Hazewinkel ideals. No compatibility of that
projection with power operations is needed.

Finally, the attaching maps, unit maps, cofiber sequences, and their
specified homotopies descend to a common set of primes, together with
the finitely many Hurewicz equalities.
At one such prime, regularity computes homology in all degrees at
once. The sphere map onto the resulting cyclic $\BP$ module fixes
its comodule structure. This last step is what produces an ordinary
finite Smith--Toda complex, rather than only an ultraproduct or
chromatically localized realization.

Three different forms of finite control enter this proof. Rationality
is tested one fixed homotopy degree at a time. The estimate at weights
growing with $p$ uses a finite-cell comparison of whole primewise
modules. The final realization retains only a finite diagram of
attaching maps and homotopies. Each has its own justification; no
step requires an intersection of infinitely many large prime sets.

\subsection{Organization}
Section~\ref{sec:foundations} constructs the rational graded models
and their finite realizations. Sections~\ref{sec:power}
and~\ref{sec:layer} establish the uniform permutation power and its
Euler quotient. Sections~\ref{sec:universal}
and~\ref{sec:connectivity} give universal coefficient algebras and
the large-weight estimate. Section~\ref{sec:normalization}
normalizes the power and constructs the successive attaching classes.
Section~\ref{sec:cobordism} proves the coefficient identity needed
for detection, and Section~\ref{sec:realization} completes the
Hurewicz induction and finite descent.
